% ============================================================
\chapter{Stochastic Processes --- Definitions}
\label{ch:definitions}
% ============================================================

What exactly is a stochastic process? The question admits two equally valid answers, and the tension between them runs through the entire theory. On one hand, a stochastic process is a family of random variables indexed by time: at each instant $t$, we draw a value from some distribution. On the other hand, it is a single random function---a trajectory, a path, a realisation chosen at random from a space of possible histories. The first viewpoint is algebraic and distributional; the second is geometric and pathwise.

This dual perspective was clarified in the early twentieth century by Andrei Kolmogorov and Joseph Doob, who built the rigorous foundations of stochastic process theory. Kolmogorov's consistency theorem (1933) showed that specifying all finite-dimensional distributions suffices to construct the process, while Doob's work on regularity conditions revealed when one can demand that the paths be continuous or at least well-behaved. Understanding this duality---distributions versus paths---is the first step toward mastering the theory.

\section{Definition of a Stochastic Process}

\begin{definition}[Stochastic Process]
Let $(\Omega, \mathcal{F}, \PP)$ be a probability space and $T$ an index set
(typically $T = \N$, $\Z_+$, or $\R_+$). A \emph{stochastic process} with values
in a measurable space $(E, \mathcal{E})$ is a family of random variables
\[
  X = (X_t)_{t \in T}, \quad X_t : \Omega \to E,
\]
where each $X_t$ is $\mathcal{F}/\mathcal{E}$-measurable.

Equivalently, $X$ is a measurable map
$X : (T \times \Omega, \mathcal{B}(T) \otimes \mathcal{F}) \to (E, \mathcal{E})$
(when joint measurability is required).
\end{definition}

\begin{definition}[Sample Path]
For a fixed $\omega \in \Omega$, the map $t \mapsto X_t(\omega)$ is called the
\emph{sample path} (or \emph{trajectory}, or \emph{realisation}) of the process
associated with $\omega$.
\end{definition}

\begin{example}
\begin{enumerate}
  \item \textbf{Simple random walk}: let $(Y_n)_{n \geq 1}$ be i.i.d.\ with
    $\PP(Y_n = 1) = \PP(Y_n = -1) = 1/2$. The process $S_n = \sum_{k=1}^{n} Y_k$
    for $n \geq 0$ (with $S_0 = 0$) is a discrete-time, $\Z$-valued stochastic process.
  \item \textbf{Poisson process}: $N = (N_t)_{t \geq 0}$ with values in $\N$,
    having increasing c\`adl\`ag sample paths.
  \item \textbf{Brownian motion}: $B = (B_t)_{t \geq 0}$ with values in $\R$,
    having continuous sample paths.
\end{enumerate}
\end{example}

\section{Classification of Processes}

\begin{definition}[Discrete vs.\ Continuous Time]
\begin{itemize}
  \item A process is \emph{discrete-time} if $T$ is countable
    (e.g.\ $T = \N$ or $T = \Z$).
  \item A process is \emph{continuous-time} if $T$ is an interval in $\R$
    (e.g.\ $T = \R_+$ or $T = [0, T_0]$).
\end{itemize}
\end{definition}

\begin{definition}[State Space]
The space $(E, \mathcal{E})$ is called the \emph{state space} of the process.
\begin{itemize}
  \item If $E$ is countable (finite or infinite), we have a \emph{discrete state space}.
  \item If $E = \R^d$, we have a \emph{continuous state space}.
\end{itemize}
\end{definition}

\begin{definition}[Path Regularity]
A real-valued process $(X_t)_{t \geq 0}$ is said to be:
\begin{itemize}
  \item \emph{continuous} if a.s.\ $t \mapsto X_t(\omega)$ is continuous,
  \item \emph{c\`adl\`ag} (right-continuous with left limits) if a.s.\
    $X_t = X_{t+}$ for all $t$ and $X_{t-}$ exists,
  \item \emph{c\`agl\`ad} (left-continuous with right limits) if a.s.\
    $X_t = X_{t-}$ for all $t$ and $X_{t+}$ exists.
\end{itemize}
\end{definition}

\section{Filtrations}

\begin{definition}[Filtration]
\label{def:filtration}
A \emph{filtration} on $(\Omega, \mathcal{F}, \PP)$ is an increasing family
of sub-$\sigma$-algebras $\mathbb{F} = (\mathcal{F}_t)_{t \in T}$:
\[
  s \leq t \implies \mathcal{F}_s \subset \mathcal{F}_t \subset \mathcal{F}.
\]
The space $(\Omega, \mathcal{F}, \mathbb{F}, \PP)$ is called a
\emph{filtered probability space}.
\end{definition}

\begin{intuition}[title=\textbf{Increasing Information}]
The filtration models the \emph{information available over time}. The
$\sigma$-algebra $\mathcal{F}_t$ represents the set of events that can be
observed (decided as occurred or not) at time $t$. The increase reflects the
fact that information is never lost: what is known at time $s$ remains known
at time $t > s$.
\end{intuition}

\begin{definition}[Natural Filtration]
The \emph{natural filtration} (or \emph{generated filtration}) of a process
$(X_t)_{t \in T}$ is
\[
  \mathcal{F}_t^X = \sigma(X_s : s \leq t), \quad t \in T.
\]
This is the smallest filtration with respect to which $(X_t)$ is adapted.
\end{definition}

\begin{definition}[Usual Conditions]
\label{def:usual_conditions}
A filtration $(\mathcal{F}_t)_{t \geq 0}$ satisfies the \emph{usual conditions}
(of Dellacherie--Meyer) if:
\begin{enumerate}
  \item \textbf{Completeness}: $\mathcal{F}_0$ contains all $\PP$-null sets
    of $\mathcal{F}$.
  \item \textbf{Right-continuity}: for every $t \geq 0$,
    $\mathcal{F}_t = \mathcal{F}_{t+} := \bigcap_{s > t} \mathcal{F}_s$.
\end{enumerate}
\end{definition}

\begin{remark}
Any filtration can be ``completed'' and ``regularised''. Let $\mathcal{N}$
be the collection of $\PP$-null sets. Define
\[
  \overline{\mathcal{F}}_t = \sigma(\mathcal{F}_t \cup \mathcal{N}), \quad
  \widehat{\mathcal{F}}_t = \bigcap_{s > t} \overline{\mathcal{F}}_s.
\]
The filtration $(\widehat{\mathcal{F}}_t)$ satisfies the usual conditions.
\end{remark}

\section{Adapted and Predictable Processes}

\begin{definition}[Adapted Process]
A process $(X_t)_{t \in T}$ is \emph{adapted} to the filtration
$(\mathcal{F}_t)_{t \in T}$ if for every $t \in T$, $X_t$ is
$\mathcal{F}_t$-measurable.
\end{definition}

\begin{definition}[Predictable Process --- Discrete Time]
In discrete time, a process $(X_n)_{n \geq 0}$ is \emph{predictable} with respect
to $(\mathcal{F}_n)$ if $X_0$ is deterministic (or $\mathcal{F}_0$-measurable)
and $X_n$ is $\mathcal{F}_{n-1}$-measurable for all $n \geq 1$.
\end{definition}

\begin{definition}[Predictable $\sigma$-Algebra --- Continuous Time]
\label{def:predictable_sigma}
In continuous time, the \emph{predictable $\sigma$-algebra} $\mathcal{P}$ on
$\R_+ \times \Omega$ is generated by sets of the form $(s, t] \times F$ with
$0 \leq s < t$ and $F \in \mathcal{F}_s$, and sets $\{0\} \times F_0$ with
$F_0 \in \mathcal{F}_0$.

A process is \emph{predictable} if it is measurable with respect to $\mathcal{P}$.
\end{definition}

\begin{remark}
Every adapted left-continuous process is predictable.
Every predictable process is adapted.
However, an adapted c\`adl\`ag process is not necessarily predictable.
\end{remark}

\section{Stopping Times}

\begin{definition}[Stopping Time]
\label{def:stopping_time}
Let $(\mathcal{F}_t)_{t \in T}$ be a filtration. A random variable
$\tau : \Omega \to T \cup \{+\infty\}$ is a \emph{stopping time} (or
\emph{Markov time}) if:
\begin{itemize}
  \item In discrete time ($T = \N$): $\{\tau = n\} \in \mathcal{F}_n$ for all $n \in \N$.
  \item In continuous time ($T = \R_+$): $\{\tau \leq t\} \in \mathcal{F}_t$ for all $t \geq 0$.
\end{itemize}
\end{definition}

\begin{remark}
In discrete time, $\{\tau = n\} \in \mathcal{F}_n$ is equivalent to
$\{\tau \leq n\} \in \mathcal{F}_n$. In continuous time, if the filtration is
right-continuous, the condition $\{\tau < t\} \in \mathcal{F}_t$ is equivalent.
\end{remark}

\begin{example}
\begin{enumerate}
  \item Any constant $\tau = t_0$ is a stopping time.
  \item The \emph{first passage time} $\tau_a = \inf\{t \geq 0 : X_t = a\}$
    is a stopping time if $(X_t)$ is adapted with continuous (or c\`adl\`ag under
    the usual conditions) sample paths.
  \item The \emph{first entry time} into an open set $G$:
    $\tau_G = \inf\{t \geq 0 : X_t \in G\}$ is a stopping time under the
    usual conditions.
\end{enumerate}
\end{example}

\begin{definition}[Stopped $\sigma$-Algebra]
\label{def:stopped_sigma}
If $\tau$ is a stopping time, the \emph{$\sigma$-algebra of events prior to
$\tau$} is
\[
  \mathcal{F}_{\tau} = \{A \in \mathcal{F} : \forall t \in T,\;
  A \cap \{\tau \leq t\} \in \mathcal{F}_t\}.
\]
\end{definition}

\begin{proposition}[Properties of Stopping Times]
Let $\sigma, \tau$ be stopping times.
\begin{enumerate}
  \item $\sigma \wedge \tau = \min(\sigma, \tau)$ and
    $\sigma \vee \tau = \max(\sigma, \tau)$ are stopping times.
  \item If $\sigma \leq \tau$ a.s., then $\mathcal{F}_{\sigma} \subset \mathcal{F}_{\tau}$.
  \item In discrete time, $\tau + 1$ is a stopping time.
  \item $\{\sigma \leq \tau\}$, $\{\sigma < \tau\}$, $\{\sigma = \tau\}$ belong
    to $\mathcal{F}_{\sigma} \cap \mathcal{F}_{\tau}$.
\end{enumerate}
\end{proposition}

\begin{definition}[Stopped Process]
The process \emph{stopped at} the stopping time $\tau$ is
\[
  X_t^{\tau} = X_{t \wedge \tau} = X_{\min(t, \tau)}, \quad t \in T.
\]
\end{definition}

\section{Finite-Dimensional Distributions}

\begin{definition}[Finite-Dimensional Distributions]
The \emph{finite-dimensional distributions} of a process $(X_t)_{t \in T}$ are
the laws of the random vectors $(X_{t_1}, \ldots, X_{t_n})$ for every choice
$t_1 < \cdots < t_n$ in $T$ and every $n \geq 1$.
\end{definition}

\begin{theorem}[Kolmogorov Extension Theorem]
\label{thm:kolmogorov_extension}
Let $(E, \mathcal{E})$ be a Polish space. Let $(\mu_{t_1, \ldots, t_n})$ be a
family of probability measures on $(E^n, \mathcal{E}^{\otimes n})$ satisfying
the consistency conditions:
\begin{enumerate}
  \item \textbf{Permutation invariance}: for every permutation $\pi$ of
    $\{1, \ldots, n\}$,
    \[
      \mu_{t_{\pi(1)}, \ldots, t_{\pi(n)}}(B_{\pi(1)} \times \cdots \times B_{\pi(n)})
      = \mu_{t_1, \ldots, t_n}(B_1 \times \cdots \times B_n).
    \]
  \item \textbf{Marginalisation}:
    $\mu_{t_1, \ldots, t_n}(B_1 \times \cdots \times B_n \times E) = \mu_{t_1, \ldots, t_n}(B_1 \times \cdots \times B_n)$.
\end{enumerate}
Then there exists a probability space $(\Omega, \mathcal{F}, \PP)$ and a process
$(X_t)_{t \in T}$ whose finite-dimensional distributions are the
$(\mu_{t_1, \ldots, t_n})$.
\end{theorem}

\begin{warning}[title=\textbf{Limitations of Kolmogorov's Theorem}]
Kolmogorov's extension theorem guarantees existence of a process with the correct
finite-dimensional distributions, but \emph{not} path regularity. Properties
such as path continuity require additional results (e.g.\ the Kolmogorov--\v{C}entsov
theorem).
\end{warning}

\section{Kolmogorov--\v{C}entsov Theorem}

\begin{theorem}[Kolmogorov--\v{C}entsov]
\label{thm:kolmogorov_centsov}
Let $(X_t)_{t \in [0,T]}$ be a real-valued stochastic process. If there exist
constants $\alpha > 0$, $\beta > 0$, and $C > 0$ such that
\[
  \E[\abs{X_t - X_s}^{\alpha}] \leq C \abs{t - s}^{1 + \beta}, \quad
  \forall s, t \in [0, T],
\]
then there exists a modification $(\tilde{X}_t)$ of $(X_t)$ whose sample paths
are a.s.\ H\"older continuous with exponent $\gamma$ for every $\gamma < \beta/\alpha$.
\end{theorem}

\begin{remark}
Recall that a \emph{modification} of $(X_t)$ is a process $(\tilde{X}_t)$
defined on the same probability space with $\PP(X_t = \tilde{X}_t) = 1$
for every $t$.
\end{remark}

\section{Gaussian Processes}

\begin{definition}[Gaussian Process]
A process $(X_t)_{t \in T}$ is \emph{Gaussian} if for every choice
$t_1, \ldots, t_n \in T$, the vector $(X_{t_1}, \ldots, X_{t_n})$ has a
multivariate Gaussian distribution.

A Gaussian process is entirely determined by its \emph{mean function}
$m(t) = \E[X_t]$ and its \emph{covariance function}
$K(s, t) = \Cov(X_s, X_t)$.
\end{definition}

\begin{example}
Standard Brownian motion is the Gaussian process with $m(t) = 0$ and
$K(s, t) = \min(s, t)$. The \emph{Brownian bridge} on $[0, 1]$ is the Gaussian
process with $m(t) = 0$ and $K(s, t) = \min(s, t) - st$.
\end{example}

\section{Stationarity and Ergodicity}

\begin{definition}[Strict Stationarity]
A process $(X_t)_{t \in T}$ is \emph{strictly stationary} if for every $h > 0$
and every $n \geq 1$,
\[
  (X_{t_1+h}, \ldots, X_{t_n+h}) \overset{\mathcal{L}}{=} (X_{t_1}, \ldots, X_{t_n}).
\]
\end{definition}

\begin{definition}[Weak Stationarity]
A square-integrable process $(X_t)$ is \emph{weakly stationary} (or
\emph{second-order stationary}, or \emph{wide-sense stationary}) if:
\begin{enumerate}
  \item $\E[X_t] = \mu$ is constant,
  \item $\Cov(X_s, X_t) = R(t - s)$ depends only on $t - s$.
\end{enumerate}
The function $R$ is called the \emph{autocorrelation function}.
\end{definition}

\section{Exercises}

\begin{exercise}
Show that the natural filtration of a continuous process satisfies the
right-continuity condition after completion.
\end{exercise}

\begin{exercise}
Let $\tau$ be a stopping time with respect to $(\mathcal{F}_n)_{n \geq 0}$.
Show that $X_{\tau}$ is $\mathcal{F}_{\tau}$-measurable when $(X_n)$ is adapted.
\end{exercise}

\begin{exercise}
Let $(B_t)_{t \geq 0}$ be a standard Brownian motion and $a > 0$. Show that
$\tau_a = \inf\{t \geq 0 : B_t = a\}$ is a stopping time with respect to the
completed natural filtration of $B$.
\end{exercise}

\begin{exercise}
Verify that the Brownian bridge $(X_t)_{t \in [0,1]}$ defined by
$X_t = B_t - tB_1$ is a Gaussian process and compute its covariance function.
\end{exercise}

\begin{exercise}
Let $(X_t)_{t \geq 0}$ satisfy the hypotheses of the Kolmogorov--\v{C}entsov
theorem with $\alpha = 4$ and $\beta = 1$. What H\"older exponents are possible
for the continuous modification?
\end{exercise}
